\documentclass[11pt,a4paper]{article}

% Plan corto del Grupo 29: pares y tríos contra grupos más grandes.
% Se compila directo en Overleaf (pdfLaTeX).

\usepackage[utf8]{inputenc}
\usepackage[T1]{fontenc}
\usepackage[spanish,es-noshorthands]{babel}
\usepackage{lmodern}
\usepackage[margin=2.2cm]{geometry}
\usepackage[table]{xcolor}
\usepackage{tabularx}
\usepackage{booktabs}
\usepackage{enumitem}
\usepackage{amssymb}
\usepackage[skins]{tcolorbox}
\usepackage{titlesec}
\usepackage{fancyhdr}
\usepackage{hyperref}

% ---------- Colores ----------
\definecolor{azul}{HTML}{1F4E79}
\definecolor{azulclaro}{HTML}{E8F0F8}
\definecolor{verde}{HTML}{2E7D32}
\definecolor{verdeclaro}{HTML}{EAF5EA}
\definecolor{naranjo}{HTML}{B85C00}
\definecolor{naranjoclaro}{HTML}{FFF3E0}
\definecolor{gris}{HTML}{5F6B76}

\hypersetup{colorlinks=true, linkcolor=azul, urlcolor=azul}

% ---------- Títulos de sección ----------
\titleformat{\section}{\Large\bfseries\color{azul}}{}{0pt}{}[\vspace{-0.6em}{\color{azul}\rule{\linewidth}{0.8pt}}]
\titlespacing*{\section}{0pt}{1.4em}{0.8em}

% ---------- Encabezado y pie ----------
\setlength{\headheight}{14pt}
\pagestyle{fancy}
\fancyhf{}
\lhead{\small\color{gris}Grupo 29 · Capstone ICS2122}
\rhead{\small\color{gris}Plan corto · 4 al 8 de noviembre}
\cfoot{\small\color{gris}\thepage}
\renewcommand{\headrulewidth}{0.4pt}

% ---------- Cajas ----------
\newtcolorbox{idea}[1]{enhanced, colback=azulclaro, colframe=azul, boxrule=0.8pt,
  arc=2mm, left=3mm, right=3mm, top=2mm, bottom=2mm,
  title={\bfseries #1}, coltitle=white, colbacktitle=azul}
\newtcolorbox{ojo}[1]{enhanced, colback=naranjoclaro, colframe=naranjo, boxrule=0.8pt,
  arc=2mm, left=3mm, right=3mm, top=2mm, bottom=2mm,
  title={\bfseries #1}, coltitle=white, colbacktitle=naranjo}
\newtcolorbox{meta}[1]{enhanced, colback=verdeclaro, colframe=verde, boxrule=0.8pt,
  arc=2mm, left=3mm, right=3mm, top=2mm, bottom=2mm,
  title={\bfseries #1}, coltitle=white, colbacktitle=verde}

% ---------- Checklist ----------
\newlist{tareas}{itemize}{1}
\setlist[tareas]{label=$\square$, leftmargin=1.6em, itemsep=0.25em, topsep=0.3em}
\newcommand{\opcional}{\textit{\color{gris}(si alcanza)}\ }

% Encabezado de cada día del checklist
\newcommand{\dia}[2]{%
  \vspace{0.6em}\noindent
  \colorbox{azul}{\color{white}\bfseries\strut\ #1\ }\hspace{0.6em}{\bfseries\color{azul}#2}\par\vspace{0.2em}}

\setlength{\parindent}{0pt}
\setlength{\parskip}{0.5em}
\renewcommand{\arraystretch}{1.35}

\begin{document}

% ================= TÍTULO =================
\begin{center}
  {\color{gris}\small GRUPO 29 · PLANIFICACIÓN DE BUSES ELÉCTRICOS EN LA RM}\\[0.4em]
  {\LARGE\bfseries\color{azul} Pares y tríos contra grupos más grandes}\\[0.3em]
  {\large ¿Por qué mejoramos el plan de a 2 y 3 buses, y no de a 4, 5 u 8?}\\[0.6em]
  {\color{gris} Plan corto · de hoy, miércoles 4, al \textbf{domingo 8 de noviembre}}
\end{center}

\vspace{0.5em}

% ================= LA IDEA =================
\section{La idea, en simple}

Nuestro programa trabaja en dos pasos:
\begin{enumerate}[leftmargin=1.6em, itemsep=0.2em]
  \item \textbf{Hace un plan rápido:} asigna los viajes uno a uno, en orden de hora.
  \item \textbf{Mejora el plan por grupos:} toma 2 buses (\emph{pares}) o 3 buses (\emph{tríos}), les reparte de nuevo los viajes y deja el cambio si el costo baja.
\end{enumerate}

El profesor nos pidió demostrar que esta forma de mejorar es una buena elección. La pregunta natural es \textbf{por qué grupos de 2 y 3, y no de 4, 5 u 8}. Hay dos efectos que tiran en sentidos contrarios:

\begin{idea}{El dilema que vamos a medir}
\begin{itemize}[leftmargin=1.2em, itemsep=0.2em]
  \item \textbf{Un grupo más grande} ve más combinaciones, así que en cada intento puede encontrar ahorros que un par no ve.
  \item \textbf{Pero cada intento es mucho más lento,} así que en el mismo tiempo se alcanzan a revisar muchos menos grupos.
\end{itemize}
La comparación dice cuál de los dos efectos pesa más. Si los grupos grandes no ganan en el mismo tiempo, queda demostrado que pares y tríos son la mejor elección. Si ganan, lo sabremos a tiempo para cambiar.
\end{idea}

% ================= MÉTODOS =================
\section{Contra qué nos comparamos}

\begin{center}
\rowcolors{2}{azulclaro}{white}
\begin{tabularx}{\linewidth}{>{\bfseries}p{4.2cm} X}
\toprule
\rowcolor{azul}\color{white}Método & \color{white}\textbf{Qué hace} \\
\midrule
Solo el plan rápido & No mejora nada: es el piso \\
Solo pares & Mejora en grupos de 2 buses \\
\rowcolor{verdeclaro}Pares y tríos (el nuestro) & Grupos de 2 y después grupos de 3 \\
Cuartetos & Grupos de 4 buses \\
Quintetos & Grupos de 5 buses \\
Octetos & Grupos de 8 buses \\
Mezcla creciente & Parte con pares y sube a 3, 4, 5 y 8 solo cuando el tamaño anterior ya no encuentra ahorros \\
La respuesta perfecta & Resuelve todo sin atajos. Solo funciona en casos chicos, y sirve de regla para medir a todos los demás \\
\bottomrule
\end{tabularx}
\end{center}

\textbf{Cómo se arma cada grupo.} Solo con buses que pueden intercambiar viajes: del mismo sector y con horarios que se toquen. Se prueban dos formas de elegirlos: los buses más parecidos entre sí, y al azar.

\begin{meta}{Para que la comparación sea justa}
Todos los métodos usan los mismos casos de prueba, el mismo tiempo y el mismo computador. A todos se les calcula el costo de la misma forma, y todos pasan por el mismo revisor automático, que confirma que el plan de verdad funciona.
\end{meta}

\begin{ojo}{Ojo con la licencia}
La licencia gratuita de Gurobi solo acepta problemas pequeños, y lo más probable es que un octeto se pase del límite. Hay dos salidas: usar un programa gratuito sin límite, o pedir la licencia académica de la universidad.
\end{ojo}

% ================= CHECKLIST =================
\section{Checklist de tareas}

\dia{MIÉ 4 · JUE 5}{Investigar y preparar}
\begin{tareas}
  \item \textbf{Investigar} quién más mejora planes de buses por grupos (de 2, 3 o más buses), qué tamaño de grupo usan y por qué.
  \item Averiguar cómo se demuestra en los estudios que un método es bueno: contra qué se compara y qué se mide.
  \item Revisar cada fuente (autores, año y revista) antes de citarla.
  \item Resolver la licencia: pedir la académica o instalar el programa gratuito.
  \item Crear los \textbf{casos muy chicos}: una ruta, unos 6 casos. En ellos se puede encontrar la respuesta perfecta.
  \item Crear los \textbf{casos chicos}: de 2 a 4 rutas, unos 6 casos.
\end{tareas}

\dia{VIE 6}{Programas}
\begin{tareas}
  \item Generalizar la mejora para que acepte grupos de cualquier tamaño (2, 3, 4, 5 y 8), usando la misma pieza que hoy resuelve pares y tríos.
  \item Programar las dos formas de elegir los grupos: los buses más parecidos y al azar.
  \item Programar la mezcla creciente (2 $\to$ 3 $\to$ 4 $\to$ 5 $\to$ 8).
  \item Programar la respuesta perfecta para los casos chicos.
  \item Un programa que corra todos los métodos en todos los casos y junte los resultados en una tabla.
  \item Pruebas automáticas, y confirmar que pares y tríos siguen dando exactamente lo mismo que antes.
\end{tareas}

\dia{SÁB 7}{Correr}
\begin{tareas}
  \item Casos muy chicos y chicos: todos los métodos, incluida la respuesta perfecta.
  \item Muestra de 212 rutas: todos los métodos con 5, 10 y 20 minutos, para ver quién gana con poco tiempo y quién con más.
  \item Pasar cada resultado por el revisor automático. Lo que no pasa no se usa.
  \item \opcional De noche, la red completa con el mejor grupo grande, para compararla con la nuestra.
\end{tareas}

\dia{DOM 8}{Conclusión}
\begin{tareas}
  \item Una tabla por caso y método: costo, buses, tiempo y distancia a la respuesta perfecta.
  \item Un gráfico del ahorro según el tamaño del grupo, con el mismo tiempo.
  \item Un gráfico de cuánto demora un intento según el tamaño del grupo.
  \item Comprobar si las diferencias son reales o casualidad.
  \item Una página con la conclusión y lo que dice la investigación, lista para el informe y para una lámina.
\end{tareas}

% ================= CRITERIO =================
\section{Qué tiene que mostrar el domingo}

\begin{meta}{Podemos decir «pares y tríos es la mejor elección» si se cumple todo esto}
\begin{enumerate}[leftmargin=1.4em, itemsep=0.3em]
  \item Con el mismo tiempo, ni los cuartetos, ni los quintetos, ni los octetos, ni la mezcla creciente son más de \textbf{1\,\%} más baratos que el nuestro.
  \item En los casos chicos quedamos cerca de la respuesta perfecta, idealmente a menos de \textbf{2\,\%}.
  \item Los tríos mejoran lo que dejan los pares.
\end{enumerate}
\end{meta}

Si algo de esto no se cumple, el domingo decidimos con los datos en la mano si cambiamos el método. Una comparación honesta puede salir a favor o en contra, y saberlo antes del informe es justamente lo que buscamos.

\end{document}
